\section*{Chapitre \ref{ch:g12:continental-crust} --- La croûte continentale}

\begin{solution}{exo:g12:continental-crust:1}
\omterm{def:g12:continental-crust:continental-crust}{Croûte continentale}~: de 30 à \qty{70}{km}, environ \qty{41}{km} en
moyenne~; roches granitiques, gneiss et autres roches métamorphiques et
\omterm{def:g11:igneous-rocks:volcanic-plutonic}{plutoniques} sous une couverture sédimentaire~; environ \qty{61}{\percent}
de silice en moyenne~; environ \qty{2830}{kg/m^3}~; jusqu'à
4.03~milliards d'années (4.4 pour des grains). \omterm{def:g11:ocean-ridges:oceanic-lithosphere}{Croûte océanique}~: environ
\qty{7}{km}~; basalte et \omterm{def:g11:ocean-ridges:gabbro}{gabbro}~; environ \qty{50}{\percent} de silice~;
environ \qty{3000}{kg/m^3}~; moins de 200 millions d'années.
\end{solution}

\begin{solution}{exo:g12:continental-crust:2}
\qty{60.6}{\percent}~: \omterm{def:g11:igneous-rocks:felsic-mafic}{intermédiaire}, comme une \omterm{def:g11:igneous-rocks:rock-names}{diorite}.
\qty{66.6}{\percent}~: \omterm{def:g11:igneous-rocks:felsic-mafic}{felsique}, comme un granite.
\end{solution}

\begin{solution}{exo:g12:continental-crust:3}
La \omterm{def:g11:ocean-ridges:oceanic-lithosphere}{croûte océanique}, environ \qty{3000}{kg/m^3}, repose sur le manteau
lithosphérique, et la vieille \omterm{def:g11:ocean-ridges:oceanic-lithosphere}{lithosphère océanique} froide est, dans son
ensemble, assez dense pour s'enfoncer. La \omterm{def:g12:continental-crust:continental-crust}{croûte continentale}, environ
\qty{2830}{kg/m^3}, bien plus légère que le manteau (environ
\qty{3380}{kg/m^3}), flotte trop bien~: un continent qui arrive à une fosse
arrête la \omterm{def:g11:subduction-zones:subduction}{subduction}.
\end{solution}

\begin{solution}{exo:g12:continental-crust:4}
(1) Un cœur ancien et stable d'un continent, qui ne se déforme plus.
(2) Plus d'un milliard d'années. (3) Des roches métamorphiques et
\omterm{def:g11:igneous-rocks:volcanic-plutonic}{plutoniques}, les \omterm{def:g12:continental-crust:root}{racines} profondes d'anciennes montagnes.
\end{solution}

\begin{solution}{exo:g12:continental-crust:5}
$4030/190 \approx 21$ fois. Le fond océanique est fabriqué aux dorsales et
détruit aux \omterm{def:g11:subduction-zones:subduction}{zones de subduction}~: aucune partie ne survit beaucoup plus de
200 millions d'années.
\end{solution}

\begin{solution}{exo:g12:continental-crust:6}
(1) $0.306 \times 210.4 \approx \num{64.4}\text{ millions de }\unit{km^2}$.
(2) $210.4 - 64.4 = \num{146.0}\text{ millions de }\unit{km^2}$.
(3) $4\pi \times 6371^2 \approx \qty{5.10e8}{km^2}$~; $210.4/510.1 \approx
0.41$~: \qty{41}{\percent}.
\end{solution}

\begin{solution}{exo:g12:continental-crust:7}
(1) $41100 \times 2830 \approx \qty{1.16e8}{kg}$. (2) $7100 \times 3000 =
\qty{2.13e7}{kg}$. (3) Environ 5.5 fois.
\end{solution}

\begin{solution}{exo:g12:continental-crust:8}
\omterm{def:g12:continental-crust:root}{Racine}~: $2.5 \times 2830/550 \approx \qty{12.9}{km}$. La croûte est plus
épaisse de $2.5 \times 3380/550 \approx \qty{15.4}{km}$~: environ
\qty{50}{km} en tout, dans l'intervalle de 50 à \qty{70}{km} mesuré sous
les chaînes de montagnes.
\end{solution}

\begin{solution}{exo:g12:continental-crust:9}
(1) Environ $0.103 + 0.100 + 0.102 \approx 0.30$~: \qty{30}{\percent}.
(2) La tranche la plus jeune, de 0 à 10 millions d'années (environ
\qty{10}{\percent}), presque à égalité avec celle de 20 à 30. (3) L'ancien
fond a été détruit aux \omterm{def:g11:subduction-zones:subduction}{zones de subduction}~; plus le fond est ancien, plus
la part qui en a disparu est grande.
\end{solution}

\begin{solution}{exo:g12:continental-crust:10}
$120 = 0.694\,h - 0.306 \times 100$, donc $h = 150.6/0.694 \approx
\qty{217}{m}$. Ne compter que les terres émergées laisse de côté les bords
bas et immergés des continents, et la moyenne en sort bien plus élevée.
\end{solution}

\begin{solution}{exo:g12:continental-crust:11}
L'île~: un \omterm{def:g11:ocean-ridges:gabbro}{gabbro} sous des laves et un \omterm{def:g11:seismic-structure:moho}{Moho} à \qty{8}{km}, c'est une \omterm{def:g11:ocean-ridges:oceanic-lithosphere}{croûte
océanique} (l'altitude de l'île n'y change rien). La plaine~: du gneiss et un
\omterm{def:g11:seismic-structure:moho}{Moho} à \qty{38}{km}, c'est une \omterm{def:g12:continental-crust:continental-crust}{croûte continentale}.
\end{solution}

\begin{solution}{exo:g12:continental-crust:12}
(1) Exact~: $\tfrac{4}{3}\pi(6371^3 - 6329.9^3) \approx
\qty{2.0829e10}{km^3}$~; approché~: $4\pi \times 6371^2 \times 41.1
\approx \qty{2.0964e10}{km^3}$. (2) L'approximation est trop grande
d'environ \qty{0.65}{\percent}, proche de $H/R \approx 0.0065$~: pour une
couche mince, l'erreur est négligeable.
\end{solution}

\begin{solution}{exo:g12:continental-crust:13}
(1) $60.6 = (66.6 + 2x)/3$, donc $x = (181.8 - 66.6)/2 = \qty{57.6}{\percent}$.
(2) \omterm{def:g11:igneous-rocks:felsic-mafic}{Intermédiaire}. (3) La \omterm{def:g12:continental-crust:continental-crust}{croûte continentale} profonde est moins riche en
silice que la croûte supérieure~: davantage de minéraux \omterm{def:g11:igneous-rocks:felsic-mafic}{mafiques}, plus
riches en fer, en magnésium et en calcium.
\end{solution}

\begin{solution}{exo:g12:continental-crust:14}
(1) $3 \times 10^8/3 = 10^8$ ans, soit 100 millions d'années. (2) Le même
ordre de grandeur que les âges du fond océanique~: les trois quarts ont
moins de 100 millions d'années, aucun plus de 190 environ. (3) Il doit être
détruit aussi vite qu'il est fabriqué~: il replonge dans le manteau aux
\omterm{def:g11:subduction-zones:subduction}{zones de subduction}.
\end{solution}

\begin{solution}{exo:g12:continental-crust:15}
(1) $8.6 \times 10^9 / 4.0 \times 10^9 \approx \qty{2.2}{km^3}$ par an.
(2) Un rythme constant n'est qu'une possibilité~: la croûte a pu croître
vite au début et être recyclée depuis, ou croître tard. Les roches et les
grains les plus anciens prouvent qu'une croûte existait tôt~; ils ne
peuvent pas dire combien, car une croûte ancienne peut être érodée ou
renvoyée dans le manteau et ne laisser que peu de traces.
\end{solution}

\begin{solution}{pb:g12:continental-crust:1}
\textbf{1.} $4\pi \times 6371^2 \approx \qty{5.10e8}{km^2}$.

\textbf{2.} $210.4/510.1 \approx 0.41$~: \qty{41}{\percent}.

\textbf{3.} $0.306 \times 210.4 \approx \num{64.4}\text{ millions de }\unit{km^2}$.

\textbf{4.} $510.1 - 210.4 \approx \num{299.7}\text{ millions de }\unit{km^2}$,
soit \qty{59}{\percent}.

\textbf{5.} Les bords des continents se trouvent sous des mers peu
profondes, sur les plateaux continentaux.

\textbf{6.} L'épaisseur, \qty{41}{km}, est très petite devant le rayon,
\qty{6371}{km}~: la courbure modifie le volume de moins de un pour cent.

\textbf{7.} $210.4 \times 10^6 \times 41.1 \approx \qty{8.65e9}{km^3}$.

\textbf{8.} \qty{8.65e18}{m^3}.

\textbf{9.} $210.4 \times 10^6 \times 36 \approx \qty{7.57e9}{km^3}$.

\textbf{10.} $(8.65 - 7.57)/7.57 \approx 0.14$~: \qty{14}{\percent}.

\textbf{11.} $8.65 \times 10^{18} \times 2830 \approx \qty{2.45e22}{kg}$.

\textbf{12.} $41100 \times 2830 \approx \qty{1.16e8}{kg}$.

\textbf{13.} $299.7 \times 10^6 \times 7.1 \approx \qty{2.13e9}{km^3}$,
soit \qty{2.13e18}{m^3}~; multiplié par 3000~: environ \qty{6.4e21}{kg}.

\textbf{14.} La \omterm{def:g12:continental-crust:continental-crust}{croûte continentale}, environ 3.8 fois plus lourde~: elle est
presque six fois plus épaisse, ce qui l'emporte sur sa surface et sa masse
volumique plus faibles.

\textbf{15.} Elle est faite de roches granitiques, riches en quartz et en
feldspaths, les plus légers des minéraux constitutifs des roches~; la
\omterm{def:g11:ocean-ridges:oceanic-lithosphere}{croûte océanique} est faite de basalte et de \omterm{def:g11:ocean-ridges:gabbro}{gabbro}, riches en pyroxène et
en olivine.

\textbf{16.} $2.45 \times 10^{22}/5.97 \times 10^{24} \approx 0.0041$.

\textbf{17.} $3 \times 0.41 \times (41.1/6371) \times (2830/5513) \approx
0.0041$~: le même résultat.

\textbf{18.} $6.4 \times 10^{21}/5.97 \times 10^{24} \approx 0.0011$~:
environ un millième.

\textbf{19.} Environ \qty{0.4}{\percent}.

\textbf{20.} La \omterm{def:g12:continental-crust:continental-crust}{croûte continentale} représente environ
\textbf{\qty{0.4}{\percent} de la masse de la Terre}.
\end{solution}
